\section{Methods}\label{sec:methods}

\subsection{Calibration data}\label{sec:calib}

IBM distributes snapshots of device calibrations with the fake backends of qiskit-ibm-runtime
\cite{qiskit_ibm_runtime} (version 0.50.0, Apache License 2.0). We export \NumDevices{} of them to CSV files in the
column layout of IBM's calibration downloads: per qubit $T_1$, $T_2$, readout assignment error, the probability of
reading 1 after preparing 0, readout length, $\sqrt{X}$ error and gate length, reset length, and per coupler the CZ
or ECR error and gate length. The reader matches column headers by keywords, so a CSV downloaded from IBM's dashboard
loads unchanged; a qubit with any missing value is marked unusable and never placed. Table~\ref{tab:devices} lists
the devices: two Heron r3 snapshots from April 2026, and Heron r2, Heron r1 and Eagle r3 snapshots from February 2025.

\begin{table}[t]
\caption{The eight calibration snapshots: medians over usable qubits and couplers.}\label{tab:devices}
\centering
\setlength{\tabcolsep}{4pt}
% generated by analysis/summarize.py; do not edit
\begin{tabular}{llrrrrrrr}
\toprule
Device & Family & Usable & 2q & 2q error & sx error & Readout & $T_1$ & $T_2$ \\
 & (snapshot) & qubits & gate & (\%) & (\%) & error (\%) & ($\mu$s) & ($\mu$s) \\
\midrule
Aachen & Heron r3 (2026-04) & 156 & CZ & 0.15 & 0.019 & 0.74 & 232 & 255 \\
Kingston & Heron r3 (2026-04) & 155 & CZ & 0.18 & 0.023 & 0.95 & 283 & 144 \\
Marrakesh & Heron r2 (2025-02) & 156 & CZ & 0.32 & 0.023 & 0.95 & 197 & 118 \\
Fez & Heron r2 (2025-02) & 156 & CZ & 0.38 & 0.023 & 0.76 & 145 & 88 \\
Torino & Heron r1 (2025-02) & 132 & CZ & 0.41 & 0.027 & 2.29 & 185 & 141 \\
Sherbrooke & Eagle r3 (2025-02) & 126 & ECR & 0.75 & 0.024 & 1.97 & 279 & 169 \\
Brisbane & Eagle r3 (2025-02) & 127 & ECR & 0.77 & 0.024 & 2.00 & 232 & 150 \\
Kyiv & Eagle r3 (2025-02) & 127 & ECR & 1.15 & 0.025 & 1.27 & 287 & 118 \\
\bottomrule
\end{tabular}

\end{table}

\subsection{Calibrated noise on NVIDIA's memory circuit}\label{sec:noise}

The experiment is a rotated surface code memory of distance $d$ over $d$ rounds in the X or Z basis, built by
NVIDIA's \texttt{MemoryCircuit} with boundary detectors \cite{nvidia_ising_code}: the circuit on which the pre-decoder
was trained and evaluated. NVIDIA's noise model places noise in every round except the last, which is perfect, and
reads the data qubits without noise. We generate the circuit with distinct placeholder probabilities, one per class
of fault location, and replace every placeholder by a channel computed from the calibration of the device qubit or
coupler assigned to that location (Table~\ref{tab:noise}). Detector and observable definitions are untouched, so the
network input, the kernel input and the detector order seen by PyMatching are those of NVIDIA's pipeline.

\begin{table}[t]
\caption{Fault locations of NVIDIA's circuit and their calibrated channels. $r$ is an error per gate from randomized
benchmarking; idle channels use the duration of the layer in which the qubit waits.}\label{tab:noise}
\centering
\begin{tabular}{p{3.2cm}p{8.3cm}}
\toprule
Location & Calibrated channel \\
\midrule
Two qubit gate & two qubit depolarizing with $5r_{2q}/4$ on the pair; one qubit depolarizing with $3r_{sx}$ on the
target (the two Hadamards that turn CZ or ECR into CNOT) \\
Idle during gates & Pauli twirled $T_1$ and $T_2$ decay over the longest coupler of the layer plus two $\sqrt{X}$ gates \\
Idle during readout & the same, over the longest ancilla readout \\
Idle during reset & the same, over the longest ancilla reset \\
Measurement & bit flip with the readout assignment error \\
Preparation & bit flip with half of the probability of reading 1 after preparing 0 \\
\bottomrule
\end{tabular}
\end{table}

For an idle of duration $t$ the Pauli twirl of amplitude and phase damping \cite{ghosh2012decoherence} gives
\begin{equation}
p_x = p_y = \frac{1 - e^{-t/T_1}}{4}, \qquad p_z = \frac{1 - e^{-t/T_2}}{2} - p_x, \label{eq:idle}
\end{equation}
with $T_2$ clipped to $2T_1$. Average gate errors $r$ are converted to depolarizing probabilities with $3r/2$ for one
qubit and $5r/4$ for two qubits. Measurement and preparation flips are told apart by the next operation on the qubit:
a flip followed by a measurement of that qubit is a readout error, any other is a preparation error. The preparation
error takes half of IBM's ``Prob meas1 prep0'' because that quantity contains both the preparation and the readout of
a prepared 0. An error scale $s$ multiplies every probability and every decay rate ($t \to st$ in
Eq.~\ref{eq:idle}), so $s = 0.5$ models a device with half of every error.

\textbf{Placement.} IBM devices are heavy hexagonal and the rotated surface code needs a square grid. We do not embed
the code in the device graph, which would require flag qubits or SWAPs \cite{chamberland2020heavyhex}. Instead, each
code qubit gets a distinct device qubit and each code coupler a distinct device coupler, in one of three modes:
\emph{median} (every location gets the device median, a homogeneous circuit), \emph{best} (the lowest error usable
qubits and couplers) and \emph{screened} (a seeded random draw after removing the worst 10\% of every metric). The
noise is therefore drawn from each device's measured distribution, not a model of running the code on that chip.

\textbf{Experimental points.} Every device is run in the median and best modes at $s = 1$, $0.5$ and $0.25$ and
$d = 3$ and $5$; the two Heron r3 devices are also run with three screened draws at $s = 0.5$. Each point is sampled
with Stim \cite{gidney2021stim}, 200,000 shots per basis, with fixed seeds that make every decoder see the same shots.
The X and Z experiments are pooled, giving \NumPoints{} points of 400,000 shots.

\subsection{Decoders}\label{sec:decoders}

All matching decoders use PyMatching 2.4 \cite{higgott2022pymatching,higgott2025sparse} on the decomposed detector
error model of a circuit:
\begin{description}
\item[PM] plain matching on the calibrated circuit, so every edge has the weight of its own qubits' rates;
\item[PM median] plain matching on the median circuit of the same device, which knows only the averages;
\item[PM-corr] matching with correlations enabled, on the calibrated circuit;
\item[Pre-decoder + PM] NVIDIA's Fast network (four 3D convolutions of width 128, receptive field 9, trained at
$d = 9$ on uniform noise \cite{nvidia_fast_card}) in float32 or as the INT8 integer model, followed by PM on the
residual syndrome, through NVIDIA's evaluation module;
\item[Pre-decoder + PM-corr] the float network followed by PM-corr.
\end{description}
The integer model uses INT8 weights with one scale per output channel, INT8 activations with per layer scales taken
at the 99.99th percentile of 1,024 windows sampled from the same circuit with a different seed, and the GELU as a
lookup table; it is the model the FPGA kernel computes
\cite{ali2026nodes}. A logical failure is a shot whose predicted observable flip differs from the true one. For two
decoders $a$ and $b$ on the same shots we report the ratio of failures and the two sided exact sign test on the shots
where exactly one of them fails; a difference is called significant at $p < 0.05$. The logical error per round is
$\epsilon = \tfrac12\bigl(1 - (1 - 2P)^{1/r}\bigr)$ for a failure probability $P$ over $r$ rounds.

\subsection{One U55C card}\label{sec:card}

The card runs the bitstream of the earlier four card study \cite{ali2026nodes} (SHA256 prefix d399f8d0) on one Alveo
U55C \cite{amd_ds978}: two instances of an integer kernel whose window size, up to $13 \times 13 \times 13$, and
weights are launch arguments. A C++ host on XRT \cite{amd_xrt} streams the test windows of a point to both instances
and reads back one decision bit per site. On the first 4,096 windows of every run, both instances are checked bit
exact against the integer model. The decisions of all 200,000 windows are then passed, as fixed network outputs,
through the same evaluation module and PM, and the logical failures are compared shot by shot with the CPU integer
model. Throughput is measured with both instances busy and two launches in flight, 64 windows per launch, over three
repetitions of about ten seconds; latency is the mean time of a one window launch with one launch in flight, PCIe
transfers and runtime included.
For comparison, the float network is timed in PyTorch on the \CpuName{} host with \CpuThreads{} threads, network
forward pass only, in batches of 4,096 windows.

\subsection{Teaching example}\label{sec:teaching}

Distances 1 and 2 cannot correct errors and are not part of the evaluation. The repository includes a CPU example
that runs them, together with $d = 3$ and $5$, from any calibration CSV with the device medians: $d = 1$ is one
unencoded qubit kept for the duration of the rounds, $d = 2$ discards every shot with a detection event, and $d = 3$
and $5$ are decoded by PM. With the Kingston medians over five rounds it gives logical errors of \TeachDOneOne\% for
the bare qubit, \TeachDTwoKeptOne\% for the kept shots of $d = 2$ (with \TeachDTwoDiscardOne\% of shots discarded),
\TeachDThreeOne\% for $d = 3$ and \TeachDFiveOne\% for $d = 5$ at $s = 1$, and \TeachDFiveQuarter\% for $d = 5$ at
$s = 0.25$. It uses Stim's generated circuits with a simpler noise model and is not compared with the results below.
